\documentclass[12pt,a4paper,twoside]{report}

\usepackage[margin=2cm]{geometry}
\usepackage{graphicx}
\usepackage{amsmath}
\usepackage{amssymb}
\usepackage{hyperref}
\usepackage[numbers,round,sort&compress]{natbib}
\usepackage{fancyhdr}
\usepackage{setspace}

\onehalfspacing

\hypersetup{
    colorlinks=false,
    pdfborder={0 0 0}
}

\pagestyle{fancy}
\fancyhf{}
\fancyhead[LE,RO]{\thepage}
\fancyhead[LO]{\slshape \rightmark}
\fancyhead[RE]{\slshape \leftmark}

\newcommand{\reporttitle}{Investigating the Impact of Category Prevalence and Variance Scaling on LASSO Variable Selection}
\newcommand{\reportauthor}{Hana Wang}
\newcommand{\cid}{Add CID here}
\newcommand{\degreetype}{Health Data Analytics and Machine Learning}
\newcommand{\submissiondate}{July 2026}
\newcommand{\wordcount}{Add word count here}

\begin{document}

\begin{titlepage}
    \centering

    {\Large IMPERIAL COLLEGE LONDON\par}
    \vspace{0.5cm}
    {\Large SCHOOL OF PUBLIC HEALTH\par}

    \vspace{3cm}

    {\LARGE\bfseries Investigating the Impact of Category Prevalence\par}
    \vspace{0.3cm}
    {\LARGE\bfseries and Variance Scaling on LASSO Variable Selection\par}

    \vspace{2cm}

    {\large \reportauthor\par}
    \vspace{0.4cm}
    {\large CID: \cid\par}
    \vspace{0.4cm}
    {\large Degree Type: \degreetype\par}
    \vspace{0.4cm}
    {\large Date of Submission: \submissiondate\par}
    \vspace{0.4cm}
    {\large Word Count: \wordcount\par}

    \vfill

    {\large Submitted in partial fulfilment of the requirements for the MSc degree in Health\par}
    {\large Data Analytics and Machine Learning of Imperial College London\par}
\end{titlepage}

\pagenumbering{roman}
\setcounter{page}{1}

\chapter*{Abstract}
\addcontentsline{toc}{chapter}{Abstract}

\textbf{Background:} Category prevalence may affect the variance of encoded categorical predictors, which can influence LASSO-based variable selection.

\textbf{Aims/Objectives:} This project investigates whether category rarity and variance scaling affect variable selection performance in standard LASSO and stability-selection LASSO.

\textbf{Methods:} Simulation studies will compare different scaling methods, rarity levels, predictor correlations, and noise settings. Performance will be assessed using F1 score, recall, and precision.

\textbf{Results:} The analysis will examine whether rare categories, scaling choices, correlation, and noise change selection performance, and whether stability selection reduces sensitivity to these factors.

\textbf{Conclusions:} The findings will help clarify how categorical predictor prevalence and scaling should be considered when applying LASSO-based variable selection.

\cleardoublepage

\chapter*{Acknowledgements}
\addcontentsline{toc}{chapter}{Acknowledgements}

Add acknowledgements here. Remove names from the anonymous assessed copy.

\cleardoublepage

\tableofcontents
\cleardoublepage
\listoftables
\addcontentsline{toc}{chapter}{List of Tables}
\cleardoublepage
\listoffigures
\addcontentsline{toc}{chapter}{List of Figures}
\cleardoublepage

\pagenumbering{arabic}
\setcounter{page}{1}

\chapter{Analysis Plan}

\section{Overall Aim}

\begin{itemize}
  \item Investigate whether category prevalence and variance scaling affect LASSO-based variable selection.
  \item Compare standard LASSO and stability-selection LASSO.
  \item Evaluate performance using F1 score, recall, and precision.
  \item Compare results across rarity, correlation, and noise settings.
\end{itemize}

\section{Research Question 1: Fully Categorical Predictors}

\begin{itemize}
  \item Without scaling, how does category rarity affect variable-selection performance across different scenarios?
  \item Does z-score scaling improve variable-selection performance relative to no scaling across different scenarios?
  \item Examine whether scaling and rarity affect selection when all predictors are categorical, represented as binary predictors for simplicity.
  \item Start from a balanced all-binary baseline, then introduce increasingly rare categories, signal strength variation, and correlation.
  \item For LASSO, compare $\lambda_{\min}$ and $\lambda_{1se}$.
  \item For stability-selection LASSO, compare $n_{\mathrm{cat}} = 0$ and $n_{\mathrm{cat}} = 3$.
  \item Apply two scaling methods: no scaling and z-score scaling.
\end{itemize}

\section{Research Question 2: Mixed Predictors}

\begin{itemize}
  \item In datasets containing both binary and continuous predictors, which scaling strategy produces the best variable-selection performance across different scenarios?
  \item Repeat the main rarity analysis with 50\% binary predictors and 50\% continuous predictors.
  \item Apply three scaling methods: no scaling, z-score scaling of all predictors, and two-standard-deviation scaling of continuous predictors only.
\end{itemize}

\section{Simulation Scenarios for Both Research Questions}

\begin{itemize}
  \item \textbf{Scenario 0 for Research Question 1:} baseline setting with no rare categories. This acts as a sanity check and should be similar regardless of rarity and scaling.
  \item \textbf{Scenario 1:} increasing rarity under lower-noise conditions, with signal strength around $ev_{xy} = 0.5$.
  \item \textbf{Scenario 2:} increasing rarity under medium-noise conditions, with weaker signal such as $ev_{xy} = 0.2$ or the closest available simulated value.
  \item \textbf{Scenario 3:} increasing rarity under higher-noise conditions, with weaker signal such as $ev_{xy} = 0.05$ or the closest available simulated value.
  \item Vary predictor correlation across no correlation, low correlation, and high correlation settings.
  \item Repeat all scenarios 100 times to obtain empirical 95\% percentile intervals across simulation seeds.
\end{itemize}

\section{Main Assessment and Plotting Strategy}

\begin{itemize}
  \item Produce separate small-multiple figures for LASSO and stability-selection LASSO.
  \item Use columns for F1 score, recall, and precision.
  \item Use rows for no correlation, low correlation, and high correlation.
  \item Divide each figure into lower-, medium-, and higher-noise blocks.
  \item Use rarity severity on the x-axis within each small plot.
  \item In the LASSO figure, each panel will contain lines depending on $\lambda$ tuning settings crossed with scaling methods.
  \item In the stability-selection LASSO figure, each panel will contain lines representing the $n_{\mathrm{cat}}$ settings crossed with the scaling methods.
  \item Generate plots for overall variable-selection performance and, where relevant, binary-predictor selection, rare-category binary-predictor selection, and continuous-predictor selection.
\end{itemize}

\chapter{Introduction}

\begin{itemize}
  \item Introduce LASSO as a commonly used method for variable selection in high-dimensional data.
  \item Explain that categorical predictors can create challenges because category prevalence affects predictor variance.
  \item Introduce the concern that rare categories and scaling choices may influence which variables are selected.
  \item State the main aim: to investigate the impact of category prevalence and variance scaling on LASSO-based variable selection.
\end{itemize}

\chapter{Methods}

\begin{itemize}
  \item Describe the simulation design and predictor-generation process.
  \item Explain how rarity, scaling, correlation, and noise are varied.
  \item Describe the LASSO and stability-selection LASSO settings.
  \item Define the performance metrics: F1 score, recall, and precision.
\end{itemize}

\chapter{Results}

\begin{itemize}
  \item Present the main LASSO plots across rarity, correlation, and noise settings.
  \item Present the main stability-selection LASSO plots across the same settings.
  \item Compare F1 score, recall, and precision across scaling methods.
  \item Summarise whether performance changes more strongly with rarity, scaling, correlation, or noise.
  \item Include sensitivity-analysis plots for overall, categorical, and rare-category selection performance.
\end{itemize}

\chapter{Discussion}

\begin{itemize}
  \item Interpret whether category prevalence affects variable selection performance.
  \item Discuss whether scaling improves or worsens selection under rare-category settings.
  \item Compare the behaviour of standard LASSO and stability-selection LASSO.
  \item Discuss whether the findings persist in mixed categorical and continuous data.
  \item Mention limitations of the simulation design and possible extensions.
\end{itemize}

\chapter{Conclusions}

Summarise the main conclusions of the project here.

\chapter{References}

Add references for LASSO, stability selection, scaling of binary or categorical predictors, and simulation methodology.

% Later, replace the note above with:
% \bibliographystyle{plainnat}
% \bibliography{references}

\appendix

\chapter{Supplementary Material}

Add extra simulation details, code snippets, supplementary plots, or larger tables here.

\end{document}
